\sectionguard
\section*{Source-Grounded Synthesis Across the Propers}

\subsection*{What history can and cannot say about the combination}

The textual layers do not begin as one fixed Sunday dossier. The Secret and
Postcommunion occur separately in the Veronese collection. The Old Gelasian
first supplies the checked three-oration set, but among Sunday Masses not yet
fixed to this ordinal. The Hadrianum Gregorian places the set at its Twelfth
Sunday after Pentecost. Its marginal chant cues also show that the present
chants and present orations were not then paired as they are in 1962. The
later Roman alignment is a genuine liturgical fact; it does not let an editor
project a single Carolingian or patristic compositional intention backward.

Medieval expositions preserve the same necessary stratigraphy. Honorius
Augustodunensis, \emph{Gemma animae} IV.65, interprets an older Sunday
containing the current Introit, Collect, Epistle, Gradual, Offertory, and
Communion, but with Luke~18 and another Alleluia. At IV.68 he treats Mark~7
with a different Epistle and chant field. Saint Anthony of Padua's direct
Mark~7 sermon likewise belongs to a different numbered Sunday and set. These
are valuable precedents for reading propers together, but none is evidence
that a medieval authority interpreted the exact 1962 combination.

\subsection*{Reception, fruit, and right confession}

Within that limit, the checked sources converge on a strong relation. Paul
receives and delivers a Resurrection gospel; Mark's Christ opens hearing and
makes speech right; the Alleluia frames sound as rejoicing in the true helper;
the Secret asks that active service become accepted gift and help. Gregory's
Sunday Matins lesson makes the relation explicit at one point: divine gifts
open ears, divine wisdom touches the tongue, and right speech requires action
before teaching. Bede adds catechesis and confession. Augustine and Pope
Leo~XIV qualify proclamation through the difficult secrecy command.

The result is not a theory that passive hearing is holier than action.
``Received'' and ``delivered'' are both Pauline; Mark's restored speech is a
good; the psalms praise; the Church offers. The sources instead order the
acts. Christian speech answers Christ's bodily Resurrection and the grace
that transforms a persecutor. Christian service returns what was received and
asks God for acceptance. Fruitfulness is not verbal quantity, reputation, or
self-certification.

\subsection*{One embodied economy of help}

The Gradual, Gospel, Offertory, Secret, and Postcommunion share a vocabulary
of help, bodily renewal, healing, fragility, and remedy. Ancient exegesis
makes the Resurrection link more than a modern thematic guess: Augustine and
Cassiodorus read the Gradual's flourishing flesh as rising; Greek and Latin
witnesses read Psalm~29's dedication and healing through Christ's renewed
body or renewed human nature; Irenaeus and Tertullian secure the bodily
content of Paul's gospel. Mark prevents this field from floating free of
particular suffering, while the Postcommunion includes both mind and body in
one salvation.

This synthesis is Catholic only with its limits intact. A bodily miracle does
not make disability a moral defect. Eucharistic remedy does not promise
immediate cure or eliminate ordinary medicine. Resurrection does not make the
present body irrelevant; it is precisely the body's promised end. Nor may
``whole person'' be used to collapse the distinct natural, sacramental,
pastoral, and eschatological modes in which help is given.

\subsection*{Firstfruits within the holy house}

Psalm~67's house of concord and Proverbs' firstfruits belong to different
books and historical horizons. Their later reception nonetheless supplies a
bounded ecclesial relation. Hilary, Augustine, and Cassiodorus understand the
house as a people made one by grace. Bede distinguishes honestly held goods
given in alms, good works credited to God as their source, and the heavenly
storehouse; Salvian presses the economic summons more starkly as a demand not
to withhold possessions from the poor. The Collect and Secret prevent the common life from becoming a closed
exchange: divine abundance surpasses merit, and even the offered gift waits
for acceptance and medicinal fruit.

The Communion's barns and presses therefore function best as sapiential and
eschatological abundance, not as a prosperity bargain. The holy household is
recognized by concord, care for weakness, worship, truthful confession, and
the return of firstfruits. Its final wealth is the heavenly remedy in whose
fullness mind and body hope to glory.
